\documentclass[11pt,a4paper]{article}

% Paquetes esenciales
\usepackage[utf8]{inputenc}
\usepackage[spanish]{babel}
\usepackage[left=0.75in,top=0.6in,right=0.75in,bottom=0.6in]{geometry}
\usepackage{titlesec}
\usepackage{xcolor}
\usepackage{hyperref}
\usepackage{fontawesome}
\usepackage{multicol}
\usepackage{enumitem}
\usepackage{tabularx}
\usepackage{fancyhdr}
\usepackage{graphicx} % Añadido para soporte de imágenes

% Tipografía moderna para informática
\usepackage[T1]{fontenc}
\usepackage{sourcesanspro} % Fuente principal: Source Sans Pro
\usepackage{sourcecodepro} % Fuente monospace para código
\renewcommand{\familydefault}{\sfdefault} % Sans-serif por defecto

% Definir colores (estilo Deedy)
\definecolor{primary}{HTML}{2b2b2b}
\definecolor{headings}{HTML}{6A6A6A}
\definecolor{subheadings}{HTML}{333333}

% Formato de secciones (estilo Deedy)
\titleformat{\section}{\color{headings}\scshape\raggedright\large}{}{0em}{}[\color{black}\titlerule]
\titleformat{\subsection}{\color{subheadings}\raggedright\large}{}{0em}{}

% Configurar pie de página
\pagestyle{fancy}
\fancyhf{} % Limpiar headers y footers
\fancyfoot[C]{\color{headings}\footnotesize Actualizado: \today}
\renewcommand{\headrulewidth}{0pt}
\renewcommand{\footrulewidth}{0pt}

% Eliminar indentación de párrafos
\setlength{\parindent}{0pt}

% Definir comandos personalizados (estilo Deedy)
\newcommand{\runsubsection}[1]{\color{subheadings}\fontsize{12pt}{12pt}\selectfont\bfseries\uppercase{#1}\normalfont}
\newcommand{\descript}[1]{\color{subheadings}\raggedright\scshape\fontsize{11pt}{13pt}\selectfont{#1 \\}\normalfont}
\newcommand{\location}[1]{\color{headings}\raggedright\fontsize{10pt}{12pt}\selectfont{#1\\}\normalfont}
\newcommand{\sectionsep}{\vspace{8pt}}

% Configurar hyperlinks
\hypersetup{
    colorlinks=true,
    linkcolor=headings,
    urlcolor=headings,
    pdfborder={0 0 0}
}

% Definir entorno para listas compactas
\newenvironment{tightemize}{\vspace{-\topsep}\begin{itemize}\itemsep1pt \parskip0pt \parsep0pt}{\end{itemize}\vspace{-\topsep}}

\begin{document}

% CABECERA (estilo Deedy)
\begin{center}
    
    {\Huge \color{primary} \textbf{Joan Massallé Richart}} \\
    \vspace{3pt}
    {\Large \color{subheadings} \textsc{Grado en Ingeniería Informática}} \\
    \vspace{8pt}
    \color{headings}
    % \faMapMarker \space Santa Cristina d'Aro, Cataluña, 17246, España \\
    % \vspace{2pt}
    \faPhone \space +34~688~98~02~48 \space \textbullet \space
    \vspace{2pt}
    \faEnvelope \space \href{mailto:jmassalle@gmail.com}{jmassalle@gmail.com} \space \textbullet \space
    \faLinkedin \space \href{https://www.linkedin.com/in/joan-massall%C3%A9-richart-a9064a287/}{linkedin.com/in/joan-massalle-richart} \\
    % \vspace{2pt}
    % \faGlobe \space \href{http://www.elmeuteunom.com}{www.elmeuteunom.com} \\
    % \faGithub \space \href{https://github.com/elmeuteunom}{github.com/elmeuteunom}
\end{center}

\vspace{15pt}

% ESTRUCTURA DE DOS COLUMNAS (estilo Deedy)
\begin{multicols}{2}

% COLUMNA IZQUIERDA
\section{Experiencia Profesional}

\runsubsection{AdditioApp}
\descript{Data Science}
\location{Octubre 2024 – Actualidad | Girona}
\sectionsep
\begin{tightemize}
    \item Integración y limpieza de datos procedentes de diversas fuentes.
    \item Migración de datos de otras plataformas al sistema interno, garantizando la calidad e integridad de la información.
    \item Análisis avanzado de datos con Python, incluyendo el uso de librerías como Pandas y NumPy.
    \item Creación de scripts automatizados para el procesamiento y transformación de datos.
    \item Detección de patrones y tendencias.
    \item Colaboración con equipos multidisciplinarios para definir requisitos y optimizar procesos.
\end{tightemize}
\sectionsep

\runsubsection{Bonpreu Esclat}
\descript{Charcutería}
\location{Mayo 2023 – Octubre 2024 | Sant Feliu de Guíxols}

\section{Formación}

\runsubsection{Universidad de Girona}

\descript{Grado en Ingeniería Informática}
\location{2021 – 2025 | Girona}
\sectionsep
\begin{tightemize}
\item Proyecto final: Modelización y Predicción de Trayectorias Estocásticas en Sistemas Físicos con Redes Neuronales.
\item Asignaturas destacadas: Estructura y tecnología de computadores (matrícula de honor), Metodología de la Programación, Ingeniería del Software, Inteligencia Artificial, Bases de Datos.
\end{tightemize}
\sectionsep

\runsubsection{IES Ridaura}
\descript{Bachillerato Científico}
\location{2019 – 2021 | Castell-Platja d'Aro}
\sectionsep
\begin{tightemize}
\item Nota de Selectividad: 10.972
\end{tightemize}

\sectionsep

\section{Premios}

\runsubsection{10\textordmasculine{} Fórum de Trabajos de Investigación del Baix Empordà}

\descript{Plaques solars fotovoltaiques, Seguiment solar per millorar-ne el rendiment}
\location{2021 | Girona}
\sectionsep

% \section{Proyectos Destacados}

% \runsubsection{\href{https://github.com/usuario/proyecto}{Nombre del Proyecto}}
% \location{GitHub: usuario/proyecto}
% Breve descripción del proyecto, tecnologías utilizadas y resultados obtenidos.
% \sectionsep

% \runsubsection{\href{https://github.com/usuario/otro-proyecto}{Otro Proyecto}}
% \location{GitHub: usuario/otro-proyecto}
% Descripción de otro proyecto relevante para la posición.
% \sectionsep

% NUEVA COLUMNA
\columnbreak

% COLUMNA DERECHA
\section{Perfil Profesional}
Aunque me interesa cualquier ámbito relacionado con la informática, me estoy orientando al análisis y tratamiento de datos. Por eso cursaré el máster de Ciencia de Datos de la UdG durante el curso 2025-2026. Soy un apasionado de las matemáticas, la física teórica y el aprendizaje continuo. 

% Iniciaré un máster en Ciencia de Datos y busco una oportunidad de media jornada (completa en verano) para aplicar y ampliar conocimientos. 
\sectionsep

\section{Habilidades Técnicas}

\runsubsection{Lenguajes}

Python \textbullet{} Java \textbullet{} C++ \textbullet{} C \textbullet{} PHP \textbullet{} R \textbullet{} MATLAB \textbullet{} Bash \textbullet{} SQL \textbullet{} MySQL \textbullet{} Microsoft SQL \textbullet{} JSON \textbullet{} JavaScript \textbullet{} TypeScript \textbullet{} PowerShell \textbullet{} HTML \textbullet{} CSS \textbullet{} XML
\sectionsep

\runsubsection{Herramientas}

Git \textbullet{} GitHub \textbullet{} Docker \textbullet{} WSL \textbullet{} Jira \textbullet{} Slack \textbullet{} Visual Studio Code \textbullet{} IntelliJ IDEA \textbullet{} YouTrack \textbullet{} Android Studio \textbullet{} Postman \textbullet{} Bruno \textbullet{} Notion \textbullet{} Figma \textbullet{} Miro \textbullet{} Anaconda \textbullet{} Jupyter \textbullet{} VirtualBox
\sectionsep

\section{Intereses}
\textbf{Tecnología} \textbullet{} \textbf{Matemáticas} \textbullet{} \textbf{Física} \textbullet{} \textbf{Deporte} \textbullet{} \textbf{Música}

Valoro la posibilidad de trabajar a media jornada para poder compaginar trabajo y estudios.

\section{Idiomas}

\runsubsection{Catalán}
Nativo
\sectionsep

\runsubsection{Español}
Nativo
\sectionsep

\runsubsection{Inglés}
Intermedio. Nivel de 2º Bachillerato.
\sectionsep

\runsubsection{Francés}
Intermedio. Nivel DELF B1
\sectionsep

\section{Certificaciones}

\runsubsection{Cruz Roja}

\descript{Curso de primeros auxilios}
\location{2019 | Santa Cristina d'Aro}
\sectionsep

\runsubsection{Zona Jove de Santa Cristina d'Aro}

\descript{Curso de manipulación de alimentos}
\location{2019 | Santa Cristina d'Aro}
\sectionsep

% React Course (¡CUANDO LO TERMINE!)

\end{multicols}

\end{document}
\end{document}
